\chapter{The Agent Ecosystem}\index{Agent Ecosystem}\label{ch:agent-ecosystem}

Claude Code's Task tool spawns specialized agents for different purposes. Understanding this ecosystem enables effective delegation and appropriate DCF engagement.

\section{Specialized Agent Types}

\begin{table}[H]
\centering
\begin{tabular}{lll}
\toprule
\textbf{Agent Type} & \textbf{Purpose} & \textbf{DCF Checkpoint} \\
\midrule
Explore & Codebase understanding & Review findings for completeness \\
Plan & Architecture design & Full Socratic review of plan \\
code-reviewer & Code quality analysis & Evaluate flagged issues \\
code-explorer & Deep feature analysis & Verify execution path accuracy \\
code-architect & Design blueprints & Architectural review \\
Bash & Command execution & Review commands before execution \\
\bottomrule
\end{tabular}
\caption{Agent Types and DCF Checkpoints}
\end{table}

\section{When to Use Which Agent}

\begin{itemize}
    \item \textbf{Explore (quick/medium/thorough)}: When you need to understand code you haven't seen. Use ``thorough'' for comprehensive analysis, ``quick'' for targeted searches.
    \item \textbf{Plan}: When the task requires upfront design. Plans deserve full DCF engagement before approval.
    \item \textbf{code-reviewer}: After implementing changes. Review its findings but apply judgment---not all flagged issues are real problems.
    \item \textbf{code-architect}: For new features requiring design decisions. Treat blueprints as proposals, not directives.
\end{itemize}

\section{Background Agents}\index{Background Agents}

Agents can run in the background while you continue working. This raises DCF questions about when synchronous engagement is needed.

\subsection{When to Run in Background}

\begin{itemize}
    \item Long-running explorations where you don't need immediate results
    \item Multiple independent searches that can run in parallel
    \item Tasks where you'll review output later anyway
\end{itemize}

\subsection{When to Run Synchronously}

\begin{itemize}
    \item Results will immediately inform your next decision
    \item The task is high-stakes and needs real-time oversight
    \item You want to engage dialectically with the process, not just the output
\end{itemize}

\textbf{DCF Principle}: Background execution trades engagement for efficiency. Use it for exploration and information gathering; engage synchronously for judgment calls.

\section{Multi-Agent Orchestration}

Complex tasks may involve multiple agents working together:

\begin{lstlisting}
[User requests new feature]
1. Explore agent: Understand existing patterns
2. Plan agent: Design implementation approach
3. [Human reviews plan - DCF checkpoint]
4. Implementation (main agent)
5. code-reviewer agent: Quality check
6. [Human reviews findings - DCF checkpoint]
\end{lstlisting}

\textbf{DCF Application}: The orchestration itself is a form of prompt chaining. Your checkpoints should occur at decision boundaries, not between every agent.

